\section{Related Work}
\label{sec:related}
\noindent\textbf{Generator artifacts.} Most image generators build an image by up-sampling a low-resolution feature map several times. Odena \etal~\cite{odena2016checkerboard} showed that a transposed convolution whose kernel size is not divisible by its stride covers some output pixels more often than others, which produces a visible checkerboard pattern. Even when the pattern is not visible, it remains in the frequency domain. Up-sampling repeats the spectrum of the low-resolution map, so generated images have periodic peaks in their Fourier or DCT spectra that real photographs lack, and a simple classifier on the spectrum detects them reliably~\cite{zhang2019spectrum,frank2020frequency}. Each trained model leaves its own fingerprint in the noise residual of its images, which is enough to tell apart models with the same architecture~\cite{yu2019attribution,marra2019fingerprints}. That is the property our three-way task relies on. Tan \etal~\cite{tan2024npr} exploit the up-sampling directly with the neighbouring-pixel residual (NPR). Nearest-neighbour up-sampling copies one value into a $2\times2$ block, so subtracting each pixel from its neighbour in the block removes image content and leaves the up-sampling trace. This is used as an input to our network.

\noindent\textbf{Robust detection.} A detector that works on clean images can still fail once the images are post-processed. Wang \etal~\cite{wang2020cnndetect} trained a ResNet-50 on a single GAN and found that random blur and JPEG compression during training were what let it generalise to unseen generators and to compressed images. Gragnaniello \etal~\cite{gragnaniello2021critical} showed that the network's early layers matter as much. Strided convolutions and pooling at the start of a standard ResNet throw away the high-frequency detail that carries the artifacts, so removing that early down-sampling and training on unresized crops gives a large gain, a recipe Corvi \etal~\cite{corvi2023diffusion} also adopt for diffusion images. Chai \etal~\cite{chai2020patch} found that classifiers with a small receptive field, which judge each patch separately, generalise better than whole-image classifiers. This suggests the evidence is local and repeated across the image. The same down-sampling also makes CNNs sensitive to small shifts. Because a strided layer samples its input without first low-pass filtering it, moving the image by a single pixel can change the prediction~\cite{azulay2019small}. BlurPool~\cite{zhang2019blurpool} fixes this by blurring before each down-sampling step. Popescu and Farid~\cite{popescu2005resampling} showed that rescaling an image introduces periodic correlations between neighbouring pixels, the same kind of signal that up-sampling generators leave. 
